\SourceLeaf{129}%
bloß oberflächliche Aufmerksamkeit zu widmen, um besonders den praktischen Wert der zu erwartenden Vorteile beurteilen zu können.

Der Verfasser hatte seine Schrift dem französischen Nationalinstitut zur Beurteilung vorgelegt. Der Bericht, den \textsc{Delambre} darüber erstattet und den das Institut gebilligt hat, ist der Schrift selbst beigefügt. \textsc{Leonelli} war damit nicht ganz zufrieden und bringt in seinen ebenfalls beigefügten Gegenbemerkungen seine Verärgerung bisweilen sehr lebhaft zum Ausdruck. Wir werden auch auf \textsc{Delambres} Bericht und \textsc{Leonellis} Erwiderung eingehen müssen, da beide nun einen Teil des Buches selbst bilden.

Die in der \emph{ersten} Abteilung dargelegte Methode, den Logarithmus jeder vorgegebenen Größe $A$ zu berechnen, beruht auf dem Grundgedanken, diese Größe als Produkt der Form
\[
10^{\mu} a \left(1+\frac{1}{10} b\right)\left(1+\frac{1}{100} c\right)\left(1+\frac{1}{1000} d\right)\ldots
\]
darzustellen, wobei $a$, $b$, $c$, $d\ldots$ einfache ganze Zahlen bezeichnen. Erstens bilden die Logarithmen der einzelnen Faktoren eine abnehmende Folge. Man kann sich mit mehr oder weniger von ihnen begnügen, je nachdem, auf wie viele Dezimalstellen man den Logarithmus des Produkts benötigt. Zweitens kann man alle diese einzelnen Logarithmen in einer ein für alle Mal berechneten Tafel bereithalten. Eine solche Tafel liefert \textsc{Leonelli} hier zunächst für die briggischen Logarithmen auf 20 Dezimalstellen. Sie enthält zuerst die Logarithmen der ganzen Zahlen von 2 bis 9, dann die Logarithmen für alle $1+\frac{1}{10} b$, d.\,h. für 1,1; 1,2; 1,3 bis 1,9, danach für alle $1+\frac{1}{100} c$ usw. Außer den 8 ganzen Zahlen kommen also für jede Ordnung der gebrochenen Zahlen 9 Logarithmen hinzu. Weiter als bis zur 11$^{\mathrm{ten}}$ Ordnung muss die Tafel jedoch nicht ausgedehnt werden, da bei der zwölften die Logarithmen dieselben signifikanten Ziffern wie die der elften haben, und ebenso bei den folgenden. Auf diese Weise umfasst die Tafel nur 107 Logarithmen. Sie genügen, um den Logarithmus jedes Produkts der obigen Form auf 18 oder 19 Ziffern durch bloße Addition zu berechnen. Die 20$^{\mathrm{ste}}$ Stelle bleibt dabei offenbar immer unsicher. Für die hyperbolischen Logarithmen hat \textsc{Leonelli} eine ähnliche Tafel von gleichem Umfang beigefügt.

